\section{A closed-form solution for the stalling period} \label{sec:closed_form_solution}

As established in Section~\ref{sec:stalling_phenomenon}, the stalling phenomenon is fundamentally characterised by a strict periodicity in the primary iterates and a corresponding linear accumulation within the auxiliary error variables. Based on the geometric observations discussed in~\cite{DYKSTRASTALLING}, we can mathematically formalise the onset of this stalling phase.

\begin{definition}[Stalling]\label{def:stalling}
Given an iteration index $m \geq n-1$, a stalling period is defined as a sequence of iterations $i \in \{0,\dots,m_\text{stall}\}$ with $m_\text{stall} \geq n$ for which the primary iterates exhibit strict cyclical repetition, such that $w^{(m+i)} = w^{(m+i-n)}$.
\end{definition}

During a defined stalling period, the spatial coordinates of the primary variables $w^{(m+i)}$ remain constant across successive cycles. Consequently, only the auxiliary variables $e^{(m+i)}$ experience modification. As dictated by the update rule, these auxiliary variables are adjusted at each cycle by constant corrective increments $w^{(m)} - w^{(m+1)}$. 

The stalling period will therefore strictly continue until the cumulative magnitude of the auxiliary variable $e^{(m+i)}$ becomes sufficiently large to satisfy \[w^{(m+i)} + e^{(m-n+i)} \in \mathcal{H}_{[m+i]}\] for a given constraint $[m]$. Geometrically, this indicates that the specific half-space associated with the stall transitions from an active state to an inactive state. For polyhedral sets, because the auxiliary accumulation is strictly linear, the exact length of the stalling period can be deterministically pre-computed as soon as the algorithm detects the stalling condition. This computation relies exclusively on the active constraints at the onset of the stall and is formalised in Theorem~\ref{thm:nstall}.

\begin{theorem}[Length of Stalling Period]\label{thm:nstall}
Suppose that stalling commences at iteration $m$ as outlined in Definition~\ref{def:stalling}. Let the active half-spaces be denoted by the set $\mathcal{A} \coloneqq \{i \in \{0,\dots,n-1\} \mid w^{(m+i)} + e^{(m-n+i)} \notin \mathcal{H}_{[m+i]}\}$. The total length of the stalling period is given by:
\begin{align}\label{eq:nstall}
    M_\mathrm{stall} \coloneqq \min_{i \in \mathcal{S}} \left\lceil \frac{k^{(m-n+i)}}{b_{[m+i]} - (w^{(m+i)})^T a_{[m+i]}}\right\rceil,
\end{align}
where $\lceil \cdot \rceil$ represents the ceiling operator, $k^{(m-n+i)} = (e^{(m-n+i)})^T a_{[m+i]}$, and the candidate set $\mathcal{S}$ is defined as $\mathcal{S} \coloneqq \{i \in \mathcal{A} \mid (w^{(m-n+i)})^T a_{[m+i]} - b_{[m+i]} < 0\}$.
\end{theorem}

\begin{proof}
The existence of a finite integer $M_\text{stall}$ is a direct consequence of the asymptotic convergence guaranteed by the Boyle-Dykstra theorem. Following the conditional logic in the simplified polyhedral projection step~\eqref{eq:dykstra:proj:poly}, the stalling period will strictly terminate once an active half-space becomes inactive. This transition is mathematically expressed as $w^{(m+i)} + e^{(m-n+i)} \in \mathcal{H}_{[m+i]}$ for some half-space index $i \geq n$ under Definition~\ref{def:stalling}. 

As a direct consequence of the scalar update rule detailed in~\eqref{eq:km}, the only half-spaces capable of transitioning to an inactive state are those for which the spatial evaluation quantity $(w^{(m+i)})^T a_{[m+i]} - b_{[m+i]}$ is strictly negative. By Definition~\ref{def:stalling}, because the primary iterates are ``frozen'' in space, this quantity remains constant throughout the stalling duration. 

The requisite number of cycles to break the stall for any single half-space can therefore be obtained by determining the smallest possible integer $M_i$ that satisfies the inequality:
\begin{align*}
k^{(m-n+i)} + M_i \left((w^{(m+i)})^T a_{[m+i]} - b_{[m+i]}\right) < 0, \qquad i \in \mathcal{A}.
\end{align*}
Solving this inequality for $M_i$ yields $M_i = \left\lceil k^{(m-n+i)} / \left(b_{[m+i]} - (w^{(m+i)})^T a_{[m+i]}\right)\right\rceil$. The global duration of the stalling period is subsequently the minimum across all candidate half-spaces, establishing $M_\text{stall} = \min_{i \in \mathcal{S}} M_i$. The specific half-space that triggers the termination is indexed by $i_\text{stall} = \text{argmin}_{i \in \mathcal{S}} M_i$. If multiple indices $i_1 < \dots < i_j$ yield identical values such that $M_{i_1} = \dots = M_{i_j}$, the sequential update structure of Dykstra's algorithm defined in~\eqref{eq:dykstra} ensures that the half-space evaluated earliest in the sequence, $i_\text{stall} \coloneqq i_1$, is discarded first.
\end{proof}

It is important to note that the quantity $M_\text{stall}$ derived in Theorem~\ref{thm:nstall} strictly refers to the number of full $n$-step cycles of~\eqref{eq:dykstra} that must elapse. Following the conclusion of the stalling period, one specific constraint is discarded from the active half-space set, and the subset $\mathcal{A}$ must be redefined. Upon this redefinition and subsequent spatial movement, the algorithm may again navigate into a separate geometric corner that triggers a new stalling condition.